\ifdefined\gkver\else\def\gkver{student}\fi
\documentclass[\gkver]{gaokaozhenti}
\begin{document}

\chapter{1994年上海}

\begin{enumerate}
\item
%% number: 1
%% paperName: 1994年上海高考第 1 题 4 分
%% typeId: 1
%% type: 单选题
%% chapter: 其他
%% point: 物理思想
%% method: 等效替代
%% score: 4
%% degree: 850
%% duplicateId: 0
%% body:
（试点教材）求几个力的合力、计算几个电阻的总电阻、电路的简化，上述三种处理物理问题的方法都属于\xzanswer{C}
\fourchoices[answer=C]
{控制变量的方法}
{观察、实验的方法}
{等效替代的方法}
{类比的方法}
%% answer: C
%% memo:
\memoanswer{
本题原卷及材料中未提供详解，待补充。
}

\item
%% number: 1
%% paperName: 1994年上海高考第 1 题 4 分
%% typeId: 1
%% type: 单选题
%% chapter: 原子和原子核
%% point: 原子核式结构
%% method: 无
%% score: 4
%% degree: 850
%% duplicateId: 0
%% body:
提出原子核式结构模型的科学家是\xzanswer{C}
\fourchoices[answer=C]
{汤姆生}
{玻尔}
{卢瑟福}
{查德威克}
%% answer: C
%% memo:
\memoanswer{
本题原卷及材料中未提供详解，待补充。
}

\item
%% number: 2
%% paperName: 1994年上海高考第 2 题 4 分
%% typeId: 1
%% type: 单选题
%% chapter: 固体液体的性质
%% point: 固体的基本性质
%% method: 无
%% score: 4
%% degree: 850
%% duplicateId: 0
%% body:
下列说法正确的是\xzanswer{C}
\fourchoices[answer=C]
{玻璃是晶体}
{食盐是非晶体}
{云母是晶体}
{石英是非晶体}
%% answer: C
%% memo:
\memoanswer{
本题原卷及材料中未提供详解，待补充。
}

\item
%% number: 3
%% paperName: 1994年上海高考第 3 题 4 分
%% typeId: 1
%% type: 单选题
%% chapter: 运动和力
%% point: 牛顿定律的应用
%% method: 无
%% score: 4
%% degree: 650
%% duplicateId: 0
%% body:
假设汽车紧急制动后所受到的阻力的大小与汽车所受重力的大小差不多。当汽车以 $20\Ums$ 的速度行驶时，突然制动，它还能继续滑行的距离约为\xzanswer{B}
\fourchoices[answer=B]
{$40\Um$}
{$20\Um$}
{$10\Um$}
{$5\Um$}
%% answer: B
%% memo:
\memoanswer{
本题原卷及材料中未提供详解，待补充。
}

\item
%% number: 3
%% paperName: 1994年上海高考第 3 题 4 分
%% typeId: 1
%% type: 单选题
%% chapter: 功和能
%% point: 机械能守恒定律
%% method: 无
%% score: 4
%% degree: 650
%% duplicateId: 0
%% body:
（试点教材）罚点球时，足球以初速 $25\Ums$ 飞出，打在门柱上离地面 $2.4\Um$ 高处，此时足球的速度大小约为（空气阻力不计）\xzanswer{B}
\fourchoices[answer=B]
{$25\Ums$}
{$24\Ums$}
{$22\Ums$}
{$20\Ums$}
%% answer: B
%% memo:
\memoanswer{
本题原卷及材料中未提供详解，待补充。
}

\item
%% number: 4
%% paperName: 1994年上海高考第 4 题 4 分
%% typeId: 1
%% type: 单选题
%% chapter: 交流电
%% point: 变压器
%% method: 无
%% score: 4
%% degree: 550
%% duplicateId: 0
%% body:
图中理想变压器原、副线圈匝数之比 $n_{1}:n_{2} = 4:1$，原线圈两端连接光滑导轨，副线圈与电阻 $R$ 相连组成闭合回路。当直导线 AB 在匀强磁场中沿导轨匀速地向右作切割磁力线运动时，安培表 A$_{1}$ 的读数是 $12\UmA$，那么安培表 A$_{2}$ 的读数为\xzanswer{A}
\onepicture[w=6cm]{\includesvg{figs/04}}
\fourchoices[answer=A]
{0}
{$3\UmA$}
{$48\UmA$}
{与 $R$ 值大小有关}
%% answer: A
%% memo:
\memoanswer{
本题原卷及材料中未提供详解，待补充。
}

\item
%% number: 5
%% paperName: 1994年上海高考第 5 题 4 分
%% typeId: 1
%% type: 单选题
%% chapter: 电路
%% point: 动态分析
%% method: 无
%% score: 4
%% degree: 650
%% duplicateId: 0
%% body:
图示 A、B、C、D 四个电路中，电源电动势为 $E$，内阻为 $r$，定值电阻为 $R_{0}$。当滑动变阻器 $R$ 的滑动片 P 从 a 向 b 滑动时，伏特表读数将变大的电路是\xzanswer{A}
\fourchoices[answer=A, ispicture=true, h=2.6cm]
{figs/05a.png}
{figs/05b.png}
{figs/05c.png}
{figs/05d.png}
%% answer: A
%% memo:
\memoanswer{
本题原卷及材料中未提供详解，待补充。
}

\item
%% number: 5
%% paperName: 1994年上海高考第 5 题 4 分
%% typeId: 1
%% type: 单选题
%% chapter: 光学
%% point: 透镜
%% method: 无
%% score: 4
%% degree: 650
%% duplicateId: 0
%% body:
有一物体从靠近凸透镜处沿凸透镜主轴匀速移向焦点。在此过程中\xzanswer{B}
\fourchoices[answer=B]
{像匀速离开透镜}
{像加速离开透镜}
{像匀速向透镜靠近}
{像减速离开透镜}
%% answer: B
%% memo:
\memoanswer{
本题原卷及材料中未提供详解，待补充。
}

\item
%% number: 6
%% paperName: 1994年上海高考第 6 题 4 分
%% typeId: 1
%% type: 单选题
%% chapter: 电磁感应
%% point: 电磁感应现象
%% method: 无
%% score: 4
%% degree: 550
%% duplicateId: 0
%% body:
（试点教材）超导体铌三锗的临界温度约为 $18\UK$。把用铌三锗制成的圆环置于温度为 $-196\UdC$ 的液态氮中，当穿过圆环的磁通量突然增大时，圆环中\xzanswer{B}
\fourchoices[answer=B]
{将不产生感应电流}
{将产生瞬时的感应电流}
{将产生大小和方向不断变化的感应电流}
{将产生大小和方向不变的感应电流}
%% answer: B
%% memo:
\memoanswer{
本题原卷及材料中未提供详解，待补充。
}

\item
%% number: 7
%% paperName: 1994年上海高考第 7 题 4 分
%% typeId: 1
%% type: 单选题
%% chapter: 力物体的平衡
%% point: 三力平衡
%% method: 无
%% score: 4
%% degree: 550
%% duplicateId: 0
%% body:
水平横梁的一端 A 插在墙壁内，另一端装有一小滑轮 B。一轻绳的一端 C 固定于墙壁上，另一端跨过滑轮后悬挂一质量 $m = 10\Ukg$ 的重物，$\angle$CBA $= 30^{\circ}$，如图所示。则滑轮受到绳子的作用力为\xzanswer{C}
\onepicture[w=4cm]{\includesvg{figs/07}}
\fourchoices[answer=C]
{$50\UN$}
{$50\sqrt{3}\UN$}
{$100\UN$}
{$100\sqrt{3}\UN$}
%% answer: C
%% memo:
\memoanswer{
本题原卷及材料中未提供详解，待补充。
}

\item
%% number: 8
%% paperName: 1994年上海高考第 8 题 4 分
%% typeId: 1
%% type: 单选题
%% chapter: 运动和力
%% point: 动量守恒定律
%% method: 无
%% score: 4
%% degree: 450
%% duplicateId: 0
%% body:
物体 A 和 B 用轻绳相连挂在轻弹簧下静止不动，如图（a）所示。A 的质量为 $m$，B 的质量为 $M$。当连接 A、B 的绳突然断开后，物体 A 上升经某一位置时的速度大小为 $v$，这时物体 B 的下落速度大小为 $u$，如图（b）所示。在这段时间里，弹簧的弹力对物体 A 的冲量为\xzanswer{D}
\twopicture[h=4cm]
{figs/08a.png}
{figs/08b.png}
\fourchoices[answer=D]
{$mv$}
{$mv - Mu$}
{$mv + Mu$}
{$mv + mu$}
%% answer: D
%% memo:
\memoanswer{
本题原卷及材料中未提供详解，待补充。
}

\item
%% number: 9
%% paperName: 1994年上海高考第 9 题 5 分
%% typeId: 2
%% type: 多选题
%% chapter: 电场
%% point: 静电感应
%% method: 无
%% score: 5
%% degree: 550
%% duplicateId: 0
%% body:
（试点教材）绝缘着的卵形导体有一定量的电荷，达到静电平衡时，以下说法中正确的是\xzanswer{ABC}
\onepicture[w=5cm]{\includesvg{figs/09}}
\fourchoices[answer=ABC]
{电荷都在导体的外表面上，导体 a 端的电荷分布较 b 端密集}
{导体内部的电场强度等于零}
{导体外紧靠表面附近电场的方向都和导体表面垂直}
{如图，将与静电计连接的金属小球 d，从导体的 a 端移向 b 端时，静电计指针偏角逐渐减小}
%% answer: ABC
%% memo:
\memoanswer{
本题原卷及材料中未提供详解，待补充。
}

\item
%% number: 9
%% paperName: 1994年上海高考第 9 题 5 分
%% typeId: 2
%% type: 多选题
%% chapter: 光学
%% point: 光的电磁说
%% method: 无
%% score: 5
%% degree: 650
%% duplicateId: 0
%% body:
下列关于电磁波的叙述中，正确的是\xzanswer{AC}
\fourchoices[answer=AC]
{电磁波是电磁场由发生区域向远处的传播}
{电磁波在任何介质中的传播速度均为 $3.00\times10^{8}\Ums$}
{电磁波由真空进入介质传播时，波长将变短}
{电磁波不能产生干涉、衍射现象}
%% answer: AC
%% memo:
\memoanswer{
本题原卷及材料中未提供详解，待补充。
}

\item
%% number: 10
%% paperName: 1994年上海高考第 10 题 5 分
%% typeId: 2
%% type: 多选题
%% chapter: 气体性质
%% point: 气体图象
%% method: 无
%% score: 5
%% degree: 450
%% duplicateId: 0
%% body:
图中 A、B 两点表示一定质量的某种理想气体的两个状态。当气体自状态 A 变化到状态 B 时\xzanswer{ABD}
\onepicture[w=4cm]{\includesvg{figs/10}}
\fourchoices[answer=ABD]
{体积必然变大}
{有可能经过体积减小的过程}
{外界必然对气体做正功}
{气体必然从外界吸热}
%% answer: ABD
%% memo:
\memoanswer{
本题原卷及材料中未提供详解，待补充。
}

\item
%% number: 11
%% paperName: 1994年上海高考第 11 题 5 分
%% typeId: 2
%% type: 多选题
%% chapter: 运动和力
%% point: 超重失重
%% method: 无
%% score: 5
%% degree: 450
%% duplicateId: 0
%% body:
原来作匀速运动的升降机内，有一被伸长弹簧拉住的、具有一定质量的物体 A 静止在地板上，如图所示。现发现 A 突然被弹簧拉向右方。由此可判断，此时升降机的运动可能是\xzanswer{BC}
\onepicture[w=3cm]{\includesvg{figs/11}}
\fourchoices[answer=BC]
{加速上升}
{减速上升}
{加速下降}
{减速下降}
%% answer: BC
%% memo:
\memoanswer{
本题原卷及材料中未提供详解，待补充。
}

\item
%% number: 12
%% paperName: 1994年上海高考第 12 题 5 分
%% typeId: 2
%% type: 多选题
%% chapter: 磁场
%% point: 带电粒子在磁场中的运动
%% method: 无
%% score: 5
%% degree: 450
%% duplicateId: 0
%% body:
在垂直于纸面的匀强磁场中，有一原来静止的原子核。该核衰变后，放出的带电粒子和反冲核的运动轨迹分别如图中 a、b 所示。由图可以判定\xzanswer{BD}
\onepicture[w=3cm]{\includesvg{figs/12}}
\fourchoices[answer=BD]
{该核发生的是 $\alpha$ 衰变}
{该核发生的是 $\beta$ 衰变}
{磁场方向一定是垂直纸面向里}
{磁场方向向里还是向外不能判定}
%% answer: BD
%% memo:
\memoanswer{
本题原卷及材料中未提供详解，待补充。
}

\item
%% number: 13
%% paperName: 1994年上海高考第 13 题 5 分
%% typeId: 2
%% type: 多选题
%% chapter: 电磁感应
%% point: 电磁感应能量分析
%% method: 无
%% score: 5
%% degree: 450
%% duplicateId: 0
%% body:
两根光滑的金属导轨，平行放置在倾角为 $\theta$ 的斜面上，导轨的左端接有电阻 $R$，导轨自身的电阻可忽略不计。斜面处在一匀强磁场中，磁场方向垂直于斜面向上。质量为 $m$、电阻可不计的金属棒 ab，在沿着斜面、与棒垂直的恒力 $F$ 作用下沿导轨匀速上滑，并上升 $h$ 高度，如图所示。在这过程中\xzanswer{AD}
\onepicture[w=5.5cm]{\includesvg{figs/13}}
\fourchoices[answer=AD]
{作用于金属棒上的各个力的合力所作的功等于零}
{作用于金属棒上的各个力的合力所作的功等于 $mgh$ 与电阻 $R$ 上发出的焦耳热之和}
{恒力 $F$ 与安培力的合力所作的功等于零}
{恒力 $F$ 与重力的合力所作的功等于电阻 $R$ 上发出的焦耳热}
%% answer: AD
%% memo:
\memoanswer{
本题原卷及材料中未提供详解，待补充。
}

\item
%% number: 14
%% paperName: 1994年上海高考第 14 题 4 分
%% typeId: 3
%% type: 填空题
%% chapter: 原子和原子核
%% point: 半衰期
%% method: 无
%% score: 4
%% degree: 650
%% duplicateId: 0
%% body:
放射性元素铋 210 的半衰期是 5 天。$10\Ug$ 的铋 210 经过 10 天后还剩下 \tkanswer[1.6cm]{2.5}$\Ug$。
%% answer: 
\jdanswer{
2.5
}
%% memo:
\memoanswer{
本题原卷及材料中未提供详解，待补充。
}

\item
%% number: 15
%% paperName: 1994年上海高考第 15 题 4 分
%% typeId: 3
%% type: 填空题
%% chapter: 光学
%% point: 光的干涉
%% method: 无
%% score: 4
%% degree: 650
%% duplicateId: 0
%% body:
用红光做双缝干涉实验，在屏上观察到干涉条纹。在其他条件不变的情况下，改用紫光做实验，则干涉条纹间距将变 \tkanswer[1.2cm]{小}。如果改用白光做实验，在屏上将出现 \tkanswer[1.2cm]{彩} 色条纹。
%% answer: 
\jdanswer{
小，彩
}
%% memo:
\memoanswer{
本题原卷及材料中未提供详解，待补充。
}

\item
%% number: 16
%% paperName: 1994年上海高考第 16 题 4 分
%% typeId: 3
%% type: 填空题
%% chapter: 电路
%% point: 电容
%% method: 无
%% score: 4
%% degree: 650
%% duplicateId: 0
%% body:
平行板电容器充电后，两板间的电势差为 $U$。切断电源后，缩小平行板电容器两板间的距离，则电容器的电容将 \tkanswer[1.6cm]{变大}，两板间的电势差将 \tkanswer[1.6cm]{变小}。（填“不变”、“变大”或“变小”）
%% answer: 
\jdanswer{
变大，变小
}
%% memo:
\memoanswer{
本题原卷及材料中未提供详解，待补充。
}

\item
%% number: 17
%% paperName: 1994年上海高考第 17 题 4 分
%% typeId: 3
%% type: 填空题
%% chapter: 运动和力
%% point: 动量和冲量
%% method: 无
%% score: 4
%% degree: 650
%% duplicateId: 0
%% body:
据报道，今年 7 月中旬，苏梅克-列韦 9 号彗星（已分裂成若干碎块）将与木星相撞，碰撞后彗星发生巨大爆炸，并与木星融为一体。假设其中的一块质量为 $1.0\times10^{12}\Ukg$，它相对于木星的速度为 $6.0\times10^{4}\Ums$。在这块彗星与木星碰撞的过程中，它对木星的冲量是 \tkanswer[2.6cm]{$6\times10^{16}\UNs$}，损失的机械能为 \tkanswer[3.4cm]{$1.8\times10^{21}\UJ$}。（木星质量远大于彗星质量）
%% answer: 
\jdanswer{
$6\times10^{16}\UNs$，$1.8\times10^{21}\UJ$
}
%% memo:
\memoanswer{
本题原卷及材料中未提供详解，待补充。
}

\item
%% number: 18
%% paperName: 1994年上海高考第 18 题 4 分
%% typeId: 3
%% type: 填空题
%% chapter: 电路
%% point: 串并联电路
%% method: 无
%% score: 4
%% degree: 550
%% duplicateId: 0
%% body:
如图所示电路中，各电阻阻值已标出。当输入电压 $U_{AB} = 110\UV$，输出电压 $U_{CD} =$ \tkanswer[1.6cm]{$1\UV$}。
\onepicture[align=right, w=5cm]{\includesvg{figs/18}}
%% answer: 
\jdanswer{
$1\UV$
}
%% memo:
\memoanswer{
本题原卷及材料中未提供详解，待补充。
}

\item
%% number: 19
%% paperName: 1994年上海高考第 19 题 4 分
%% typeId: 3
%% type: 填空题
%% chapter: 电场
%% point: 带电粒子的偏转
%% method: 无
%% score: 4
%% degree: 450
%% duplicateId: 0
%% body:
图中虚线表示某一匀强电场区域内的若干个等势面。质子、氘核、$\alpha$ 粒子以相同的初速，沿与等势面平行的方向由 A 点进入该电场。在从上端进入电场到下端离开电场的过程中，质子、氘核、$\alpha$ 粒子的动量改变量之比是 \tkanswer[2.2cm]{$1:1:2$}，电势能改变量之比是 \tkanswer[2.2cm]{$2:1:2$}。
\onepicture[align=right, w=3.5cm]{\includesvg{figs/19}}
%% answer: 
\jdanswer{
$1:1:2$，$2:1:2$
}
%% memo:
\memoanswer{
本题原卷及材料中未提供详解，待补充。
}

\item
%% number: 20
%% paperName: 1994年上海高考第 20 题 4 分
%% typeId: 3
%% type: 填空题
%% chapter: 机械波
%% point: 波的多解性
%% method: 无
%% score: 4
%% degree: 450
%% duplicateId: 0
%% body:
一列沿 $x$ 轴正向传播的简谐波，在 $x_{1} = 10\Ucm$ 和 $x_{2} = 110\Ucm$ 处的两质点的振动图线如图所示。则质点振动的周期为 \tkanswer[1.6cm]{$4\Us$}，这列简谐波的波长为 \tkanswer[2.6cm]{$\frac{400}{4n+1}\Ucm$}（$n = 0,1,2\cdots$）。
\onepicture[align=right, w=6cm]{\includesvg{figs/20}}
%% answer: 
\jdanswer{
$4\Us$，$\frac{400}{4n+1}\Ucm$（$n = 0,1,2\cdots$）
}
%% memo:
\memoanswer{
本题原卷及材料中未提供详解，待补充。
}

\item
%% number: 20
%% paperName: 1994年上海高考第 20 题 4 分
%% typeId: 2
%% type: 多选题
%% chapter: 机械波
%% point: 波的多解性
%% method: 无
%% score: 4
%% degree: 450
%% duplicateId: 0
%% body:
（试点教材）图中实线表示 $t$ 时刻的横波波形，虚线表示经 $1\Us$ 后的波形。已知波自右向左传播，则此时质点 A 的运动方向 \tkanswer[1.8cm]{向下}。若 A、B 之间距离为 $1\Ucm$，则波速为 \tkanswer[2.8cm]{$(0.5+2n)\Ucms$}（$n = 0,1,2\cdots$）。
\onepicture[align=right, w=6cm]{\includesvg{figs/20b}}
%% answer: 
\jdanswer{
向下，$(0.5+2n)\Ucms$（$n = 0,1,2\cdots$）
}
%% memo:
\memoanswer{
本题原卷及材料中未提供详解，待补充。
}

\item
%% number: 21
%% paperName: 1994年上海高考第 21 题 4 分
%% typeId: 3
%% type: 填空题
%% chapter: 功和能
%% point: 功率
%% method: 近似与估算
%% score: 4
%% degree: 450
%% duplicateId: 0
%% body:
跳绳是一种健身运动。设某运动员的质量是 $50\Ukg$，他一分钟跳绳 180 次。假定在每次跳跃中，脚与地面的接触时间占跳跃一次所需时间的 $\frac{2}{5}$，则该运动员跳绳时克服重力做功的平均功率是 \tkanswer[1.8cm]{$75\UW$}。（$g$ 取 $10\Umsq$）
%% answer: 
\jdanswer{
$75\UW$
}
%% memo:
\memoanswer{
本题原卷及材料中未提供详解，待补充。
}

\item
%% number: 22
%% paperName: 1994年上海高考第 22 题 4 分
%% typeId: 4
%% type: 实验题
%% chapter: 光学
%% point: 光的折射
%% method: 无
%% score: 4
%% degree: 750
%% duplicateId: 0
%% body:
（单选题）以下四图表示一束白光通过三棱镜的光路图。其中正确的是\xzanswer{D}
\fourchoices[answer=D, ispicture=true, h=2.2cm]
{figs/22a.png}
{figs/22b.png}
{figs/22c.png}
{figs/22d.png}
%% answer: D
\jdanswer{
D
}
%% memo:
\memoanswer{
本题原卷及材料中未提供详解，待补充。
}

\item
%% number: 23
%% paperName: 1994年上海高考第 23 题 5 分
%% typeId: 4
%% type: 实验题
%% chapter: 电路
%% point: 电阻定律
%% method: 无
%% score: 5
%% degree: 650
%% duplicateId: 0
%% body:
（多选题）在“测定金属的电阻率”实验中，以下操作中错误的是\xzanswer{AC}
\fourchoices[answer=AC]
{用米尺量出金属丝的全长三次，算出其平均值}
{用螺旋测微器在金属丝三个不同部位各测量一次直径，算出其平均值}
{用伏安法测电阻时采用安培表内接线路，多次测量后算出其平均值}
{实验中应保持金属丝的温度不变}
%% answer: AC
\jdanswer{
AC
}
%% memo:
\memoanswer{
本题原卷及材料中未提供详解，待补充。
}

\item
%% number: 23
%% paperName: 1994年上海高考第 23 题 5 分
%% typeId: 4
%% type: 实验题
%% chapter: 力物体的平衡
%% point: 摩擦力
%% method: 无
%% score: 5
%% degree: 650
%% duplicateId: 0
%% body:
（试点教材）（多选题）在做“测定滑动摩擦系数”实验时，以下操作中错误的是\xzanswer{AB}
\onepicture[align=right, w=5cm]{\includesvg{figs/23}}
\fourchoices[answer=AB]
{只需调节滑块滑动的平面水平，不需调节定滑轮高度}
{在定滑轮一侧的小盘中加适量砝码，当滑块恰能沿平面匀速滑动时，盘中的砝码所受重力的大小就等于滑块与平面间的滑动摩擦力}
{用弹簧秤测出滑块所受重力的大小就等于滑块与平面间的正压力}
{要改变滑块与平面间的压力时，只需改变放在滑块上的砝码的质量}
%% answer: AB
\jdanswer{
AB
}
%% memo:
\memoanswer{
本题原卷及材料中未提供详解，待补充。
}

\item
%% number: 24
%% paperName: 1994年上海高考第 24 题 5 分
%% typeId: 4
%% type: 实验题
%% chapter: 力物体的平衡
%% point: 验证平行四边形实验
%% method: 无
%% score: 5
%% degree: 650
%% duplicateId: 0
%% body:
（多选题）在做“互成角度的两个力的合成”实验时，橡皮条的一端固定在木板上，用两个弹簧秤把橡皮条的另一端拉到某一确定的 O 点。以下操作中错误的是\xzanswer{ACD}
\fourchoices[answer=ACD]
{同一次实验过程中，O 点位置允许变动}
{实验中，弹簧秤必须保持与木板平行，读数时视线要正对弹簧秤刻度}
{实验中，先将其中一个弹簧秤沿某一方向拉到最大量程，然后只需调节另一弹簧秤拉力的大小和方向，把橡皮条另一端拉到 O 点}
{实验中，把橡皮条的另一端拉到 O 点时，两个弹簧秤之间夹角应取 $90^{\circ}$，以便于算出合力大小}
%% answer: ACD
\jdanswer{
ACD
}
%% memo:
\memoanswer{
本题原卷及材料中未提供详解，待补充。
}

\item
%% number: 25
%% paperName: 1994年上海高考第 25 题 5 分
%% typeId: 4
%% type: 实验题
%% chapter: 光学
%% point: 测定玻璃折射率
%% method: 无
%% score: 5
%% degree: 550
%% duplicateId: 0
%% body:
某同学在测定一厚度均匀的圆形玻璃的折射率时，先在白纸上作一与圆形玻璃同半径的圆，圆心为 O，将圆形玻璃平放在白纸上，使其边界与所画的圆重合。在玻璃一侧竖直插两枚大头针 P$_{1}$ 和 P$_{2}$，在另一侧再先后插两枚大头针 P$_{3}$ 和 P$_{4}$，使从另一侧隔着玻璃观察时，大头针 P$_{4}$、P$_{3}$ 和 P$_{2}$、P$_{1}$ 的像恰在一直线上。移去圆形玻璃和大头针后，得下图。在图中画出：
\begin{steps}
	\item
	沿 P$_{1}$、P$_{2}$ 连线方向的入射光线通过圆形玻璃后的传播方向
	\item
	光线在玻璃内的传播方向
	\item
	在光线的入射点作法线，标出入射角 $i$ 和折射角 $r$
	\item
	写出计算玻璃折射率的公式（不必计算）
\end{steps}
\onepicture[align=right, w=5cm]{\includesvg{figs/25}}
%% answer: 
\jdanswer{
\begin{enumerate}
	\item
	①②③见图
	\onepicture[w=5.5cm]{\includesvg{figs/25_an}}

	\item
	④$n = \frac{\sin i}{\sin r}$

\end{enumerate}
}
%% memo:
\memoanswer{
本题原卷及材料中未提供详解，待补充。
}

\item
%% number: 26
%% paperName: 1994年上海高考第 26 题 6 分
%% typeId: 4
%% type: 实验题
%% chapter: 电路
%% point: 电路实验
%% method: 无
%% score: 6
%% degree: 450
%% duplicateId: 0
%% body:
有一只伏特表，量程已知，内阻为 $R_{V}$。另有一电池（电动势未知，但不超过伏特表的量程，内阻可忽略）。请用这只伏特表和电池，再用一个电键和一些连接用导线，设计测量某一高值电阻 $R_{x}$ 的实验方法。（已知 $R_{x}$ 的值与 $R_{V}$ 相差不大）
\begin{steps}
	\item
	画出实验电路图。
	\item
	简要写出测量步骤和需记录的数据，导出高值电阻 $R_{x}$ 的计算式。
\end{steps}
%% answer: 
\jdanswer{
\begin{enumerate}
	\item
	①见图
	\twopicture[w=4cm]{figs/26_an1.png}{figs/26_an2.png}

	\item
	②先测出图（a）中伏特表读数 $U_{1}$，再测出图（b）中伏特表读数 $U_{2}$，根据欧姆定律：

	$\varepsilon = U_{1}$，$\varepsilon = U_{2} + \frac{U_{2}}{R_{V}}R_{x}$

	联立解得 $R_{x} = \frac{U_{1} - U_{2}}{U_{2}}R_{V}$

\end{enumerate}
}
%% memo:
\memoanswer{
本题原卷及材料中未提供详解，待补充。
}

\item
%% number: 27
%% paperName: 1994年上海高考第 27 题 10 分
%% typeId: 6
%% type: 计算题
%% chapter: 气体性质
%% point: 气体连接体
%% method: 无
%% score: 10
%% degree: 550
%% duplicateId: 0
%% body:
一圆柱形气缸直立在地面上，内有一具有质量而无摩擦的绝热活塞，把气缸分成容积相同的 A、B 两部分，如图所示。两部分气体温度相同，都是 $t_{0} = 27\UdC$。A 部分气体的压强 $p_{A0} = 1.0\times10^{5}\UPa$，B 部分气体的压强 $p_{B0} = 2.0\times10^{5}\UPa$。现对 B 部分的气体加热，使活塞上升，保持 A 部分气体温度不变，使 A 部分气体体积减小为原来的 $\frac{2}{3}$。求此时
\begin{enumerate}
	\item
	A 部分气体的压强 $p_{A}$；
	\item
	B 部分气体的温度 $T_{B}$。
\end{enumerate}
\onepicture[align=right, w=3.5cm]{\includesvg{figs/27}}
%% answer: 
\jdanswer{
\begin{enumerate}
	\item
	$p_{A} = 1.5\times10^{5}\UPa$

	\item
	$T_{B} = 500\UK$

\end{enumerate}
}
%% memo:
\memoanswer{
（1）A 部分气体作等温变化。由玻意耳定律，

$p_{A0}V = p_{A}\cdot\frac{2}{3}V$，$p_{A} = \frac{3}{2}p_{A0}$

以 $p_{A0} = 1.0\times10^{5}\UPa$ 代入得

$p_{A} = 1.5\times10^{5}\UPa$

（2）B 部分气体：初态 $p_{B0} = 2.0\times10^{5}\UPa$，$V_{B0} = V$，$T_{B0} = (273 + 27)\UK = 300\UK$

末态 $p_{B} = p_{A} + (p_{B0} - p_{A0}) = \frac{3}{2}p_{A0} + p_{A0} = 2.5\times10^{5}\UPa$

$V_{B} = V + \frac{1}{3}V = \frac{4}{3}V$

由理想气体状态方程

$\frac{p_{B0}V_{B0}}{T_{B0}} = \frac{p_{B}V_{B}}{T_{B}}$，$T_{B} = \frac{T_{B0}p_{B}V_{B}}{p_{B0}V_{B0}} = \frac{300 \times 2.5 \times 10^{5} \times \frac{4}{3}V}{2.0 \times 10^{5} \times V}\UK = 500\UK$

本题共 10 分。（1）4 分。正确列式得 2 分，结果正确得 2 分。（2）6 分。正确列式 3 分，末态 $p_{B}$ 正确 2 分，结果正确 1 分。
}

\item
%% number: 28
%% paperName: 1994年上海高考第 28 题 13 分
%% typeId: 6
%% type: 计算题
%% chapter: 功和能
%% point: 动能定理
%% method: 无
%% score: 13
%% degree: 350
%% duplicateId: 0
%% body:
如图所示，轻质长绳水平地跨在相距为 $2l$ 的两个小定滑轮 A、B 上，质量为 $m$ 的物块悬挂在绳上 O 点，O 与 A、B 两滑轮的距离相等。在轻绳两端 C、D 分别施加竖直向下的恒力 $F = mg$。先托住物块，使绳处于水平拉直状态，由静止释放物块，在物块下落过程中，保持 C、D 两端的拉力 $F$ 不变。
\begin{enumerate}
	\item
	当物块下落距离 $h$ 为多大时，物块的加速度为零？
	\item
	在物块下落上述距离的过程中，克服 C 端恒力 $F$ 做功 $W$ 为多少？
	\item
	求物块下落过程中的最大速度 $v_{m}$ 和最大距离 $H$？
\end{enumerate}
\onepicture[align=right, w=5cm]{\includesvg{figs/28}}
%% answer: 
\jdanswer{
\begin{enumerate}
	\item
	$\frac{\sqrt{3}}{3}l$

	\item
	$W = (\frac{2\sqrt{3}}{3} - 1)mgl$

	\item
	$v_{m} = \sqrt{2(2 - \sqrt{3})gl}$，$H = \frac{4}{3}l$

\end{enumerate}
}
%% memo:
\memoanswer{
\onepicture[align=right, w=5.5cm]{\includesvg{figs/28_an}}

（1）当物块所受合外力为零时，其加速度为零。此时物块下降距离是 $h$。因绳中拉力 $F$ 恒为 $mg$，所以此时悬点所受的三个拉力的方向互成夹角 $2\theta = 120^{\circ}$。从右图可知：

$h = l\tan 30^{\circ} = \frac{\sqrt{3}}{3}l$。

（2）物块下落 $h$ 时，C 端上升距离 $h' = \sqrt{h^{2} + l^{2}} - l$，克服 C 端恒力 $F$ 做功

$W = Fh' = mg(\sqrt{h^{2} + l^{2}} - l) = (\frac{2\sqrt{3}}{3} - 1)mgl$

（3）由动能定理，拉力和重力对物块做的功等于物块动能的增量。当物块下落 $x$ 时，有

$\frac{1}{2}mv^{2} = mgx - 2mg(\sqrt{x^{2} + l^{2}} - l)$ ①

物块在下落过程中先加速后减速，因此当物块加速度为零时，速度达最大值。以 $x = h = \frac{\sqrt{3}}{3}l$ 代入①，解得最大速度

$v_{m} = \sqrt{2(2 - \sqrt{3})gl}$

物块下落最大距离 $H$ 时，速度 $v = 0$，由①

$mgH = 2mg(\sqrt{H^{2} + l^{2}} - l)$ ②

$H(3H - 4l) = 0$，$H = 0$ 舍去，$H = \frac{4}{3}l$

本题共 13 分。

（1）3 分。正确得出加速度为零时两绳夹角为 $120^{\circ}$，得 2 分；正确算出下落距离 $h$ 得 1 分。

（2）4 分。正确得出 C 端上升距离 $h$ 结果 2 分；正确算出克服 C 端恒力 $F$ 做功 $W$ 得 2 分。

（3）6 分。正确列出①式 2 分，算出 $v_{m}$ 得 1 分；正确列出②式 2 分；算出物块下落最大距离 $H$ 得 1 分。
}

\item
%% number: 29
%% paperName: 1994年上海高考第 29 题 13 分
%% typeId: 6
%% type: 计算题
%% chapter: 电磁感应
%% point: 法拉第电磁感应定律
%% method: 无
%% score: 13
%% degree: 250
%% duplicateId: 0
%% body:
如图所示，两个正方形细导线框 1、2，质量都是 $m$，边长都是 $l$，每个框都在其两对角上接有短电阻丝（图中用粗黑线表示），阻值 $r_{1} = r_{1}' = r_{2} = r_{2}' = r$，其余部分电阻不计。两框叠放在水平面上，对应边相互平行，交叠点 A、C 位于所在边的中点。两框在交叠点彼此绝缘。在两框的交叠区域内存在竖直向上的匀强磁场（交叠区的导线框恰好在磁场边缘以内），磁感应强度为 $B$。设磁场在很短时间 $\Delta t$ 内均匀减小为零。不计所有摩擦。
\begin{enumerate}
	\item
	求流过电阻 $r_{1}$、$r_{2}$ 的电流 $I_{1}$、$I_{2}$ 的大小与方向。
	\item
	求磁场刚减小为零时，框 1 和 2 的速度 $v_{1}$ 和 $v_{2}$（并指明方向）。
	\item
	若两框在交叠点 A、C 不是互相绝缘，而是电接触良好，以上解答是否改变？并说明理由。
\end{enumerate}
\onepicture[align=right, w=6cm]{\includesvg{figs/29}}
%% answer: 
\jdanswer{
\begin{enumerate}
	\item
	$I_{1} = \frac{Bl^{2}}{8r\Delta t} = I_{2}$，方向均向下

	\item
	$v_{2} = v_{1} = \frac{\sqrt{2}B^{2}l^{3}}{32rm}$，方向向左

	\item
	不变

\end{enumerate}
}
%% memo:
\memoanswer{
（1）$E = \frac{\Delta\Phi}{\Delta t} = \frac{B\cdot\frac{1}{4}l^{2}}{\Delta t} = \frac{Bl^{2}}{4\Delta t}$

$I_{1} = \frac{E}{2r} = \frac{Bl^{2}}{8r\Delta t} = I_{2}$，方向均向下。

（2）在 $\Delta t$ 时间内，线框 1 中电流恒定，而磁场由 $B$ 均匀减小为零，故线框 1 的 Ar$_{1}'$ 部分所受平均安培力

$F_{Ar1'} = \frac{1}{2}BI_{1}\frac{l}{2} = \frac{Bl}{4}\cdot\frac{Bl^{2}}{8r\Delta t} = \frac{B^{2}l^{3}}{32r\Delta t}$，方向如图。

\onepicture[align=right, w=5cm]{\includesvg{figs/29_an1}}

同理，Cr$_{1}'$ 部分所受平均安培力 $F_{Cr1'} = \frac{B^{2}l^{3}}{32r\Delta t}$。

线框 1 所受合力 $F_{1} = \sqrt{2}F_{Ar1'} = \frac{\sqrt{2}B^{2}l^{3}}{32r\Delta t}$，方向向右。

同理，线框 2 所受合力 $F_{2} = \sqrt{2}F_{Ar2} = F_{1}$，方向向左。

由动量定理，$F\Delta t = mv$，线框 1 获得速度 $v_{1} = \frac{F_{1}\Delta t}{m} = \frac{\sqrt{2}B^{2}l^{3}}{32rm}$，方向向右。

同理，线框 2 获得速度 $v_{2} = v_{1} = \frac{\sqrt{2}B^{2}l^{3}}{32rm}$，方向向左。

（3）不变。

对中间小回路 Ar$_{2}$Cr$_{1}'$A 而言，感应电动势是 $E$，由对称性，在 Ar$_{2}$C 与 Cr$_{1}'$A 上的电动势都是 $\frac{E}{2}$。同理，对 Ar$_{1}$Cr$_{2}'$A 蝶形回路而言，感应电动势也是 $E$。由对称性，在 Ar$_{1}$C 与 Cr$_{2}'$A 上的电动势也都是 $\frac{E}{2}$。由此可得如图所示的等效电路。故 $r_{1}$、$r_{2}$ 上电流与原解相同，$v_{1}$、$v_{2}$ 也不变。

\onepicture[align=right, w=6cm]{\includesvg{figs/29_an2}}

本题共 13 分。

（1）4 分。正确求出电动势 2 分；正确求出 $I_{1}$、$I_{2}$ 1 分；正确标出方向 1 分。

（2）6 分。正确求出 $F_{Ar1'}$ 2 分；正确求出合力 $F_{1}$ 1 分；正确求出 $v_{1}$、$v_{2}$ 各 1 分；正确说明方向 1 分。

（3）回答不变 1 分；正确说明理由 2 分。
}

\end{enumerate}

\end{document}
